\section{Reproduction des taux de privation de l'ANSD}
\label{ann:E}

L'ANSD publie, dans son actualisation N-MODA 2024, la prévalence des
privations par indicateur et groupe d'âge sur l'EHCVM 2018/19. Le
tableau~\ref{tab:validation_ansd} confronte ces valeurs à celles obtenues avec
l'opérationnalisation retenue dans ce mémoire.

\begin{table}[H]
  \centering
  \caption{Comparaison de la prévalence des privations par indicateur : ce mémoire vs ANSD (EHCVM 2018, \%)}
  \label{tab:validation_ansd}
  \begin{threeparttable}
  \small
  \zebre
  \begin{tabular}{|l|c|c|c|c|c|c|}
    \hline
\entete Indicateur & \multicolumn{3}{|c|}{Ce mémoire} & \multicolumn{3}{|c|}{ANSD (MODA 2024)} \\ \hline
\entete & 0-4 & 5-14 & 15-17 & 0-4 & 5-14 & 15-17 \\
    \hline
    Type de toilettes                & 40,3 & 39,1 & 34,1 & 40,3 & 39,1 & 34,1 \\ \hline
    Partage des toilettes            & 20,4 & 19,4 & 18,6 & 20,4 & 19,4 & 18,6 \\ \hline
    Source d'eau non améliorée       & 11,9 & 10,0 & 7,9  & 12,0 & 10,7 & 8,5  \\ \hline
    Temps d'accès à l'eau $\geq 30$ min & 17,8 & 16,8 & 13,6 & 17,7 & 16,4 & 13,9 \\ \hline
    Débarras des ordures             & 59,3 & 58,2 & 53,0 & 59,3 & 58,2 & 53,0 \\ \hline
    Surpeuplement                    & 13,3 & 12,7 & 11,4 & 13,3 & 12,7 & 11,4 \\ \hline
    Insécurité alimentaire           & 45,0 & 46,5 & 44,9 & 45,0 & 46,5 & 44,9 \\ \hline
    Combustible solide               & 92,3 & 92,3 & 91,5 & 92,3 & 92,3 & 91,5 \\ \hline
    Accès santé (pas à pied)         & 71,0 & 69,8 & 65,4 & 79,3 & 78,5 & 75,2 \\ \hline
    Acte de naissance                & 31,8 & 30,0 & -    & 31,8 & 30,1 & -    \\ \hline
    Travail des enfants              & -    & 41,8 & -    & -    & 43,0 & -    \\ \hline
    Séparation parentale             & 36,5 & 43,0 & 52,8 & 36,5 & 41,8 & 52,8 \\ \hline
    Non-scolarisation                & -    & 45,3 & -    & -    & 45,3 & -    \\ \hline
    Illettrisme                      & -    & -    & 28,1 & -    & -    & 28,1 \\
    \hline
  \end{tabular}
  \end{threeparttable}

  {\footnotesize \textit{Source :} Calcul de l'auteur, EHCVM~; ANSD/UNICEF (MODA 2024) pour la colonne de comparaison.}
\end{table}

Une fois les estimations pondérées par les poids de sondage de l'EHCVM
(comme dans le calcul officiel), la plupart des indicateurs reproduisent les
taux de l'ANSD de façon très étroite~: plusieurs coïncident au dixième de
point (insécurité alimentaire, combustible solide, illettrisme, type de
toilettes, partage des toilettes, débarras des ordures, surpeuplement et
acte de naissance), les écarts résiduels restant inférieurs à 1,5 point. Le partage des toilettes suit ici la convention de
dénominateur de l'ANSD~: il n'est calculé que sur les ménages disposant
d'une installation privée, ceux à toilettes publiques ou sans toilette
étant exclus de cet indicateur, tout en restant privés au titre du type de
toilettes. La source d'eau non améliorée suit la définition ANSD à la lettre
(« traitement de l'eau de sources non adéquates »)~: est privé l'enfant vivant
dans un ménage dont la source de boisson est non améliorée (puits ouvert dans
la cour/concession, puits ouvert ailleurs, source non aménagée,
fleuve/rivière/lac, vendeur ambulant, autre~; à l'une ou l'autre saison) et
qui ne traite pas son eau (réponse différente de « oui » au filtre de
traitement, les « ne sait pas » et non-réponses étant comptés comme
non-traitement). Le dénominateur est plein (tous les ménages dont le type de
source est renseigné). Cette opérationnalisation reproduit étroitement les
taux ANSD (11,9~/~10,0~/~7,9 contre 12,0~/~10,7~/~8,5), l'écart résiduel
inférieur au point relevant du même calage démographique que les autres
indicateurs. Le travail des enfants ne retient que les activités posant clairement le
seuil d'au moins une heure~: du côté domestique, les cinq postes d'heures du
module emploi (courses au marché, travaux domestiques, garde d'enfants ou de
personnes âgées, corvée d'eau, corvée de bois)~; du côté économique, le
travail rémunéré ou pour propre compte (questions formulées « au moins une
heure »). L'accès à une structure de santé (module communautaire) suit la
définition ANSD à la lettre~: est privé l'enfant vivant dans une localité où
il ne peut accéder \emph{à pied} à un hôpital ou un centre de santé public,
c'est-à-dire dont le mode de transport habituel pour rejoindre la structure la
plus proche (question 2.02) n'est pas la marche. L'existence d'une structure
« sur place » n'est \emph{pas} assimilée à un accès à pied~: dans une grande
commune, un hôpital « existe dans la ville » mais s'y rejoint en voiture~; une
lecture par « OU inclusif » (existe \emph{ou} à pied) introduirait la condition
implicite « existe~$\Rightarrow$~à pied » et abaisserait artificiellement le
taux à 47,6\,\%. La lecture stricte donne 71,0~/~69,8~/~65,4, proche de l'ANSD
(79,3~/~78,5~/~75,2)~; l'écart résiduel tient aux localités disposant d'une
structure sur place, pour lesquelles le mode d'accès n'est pas renseigné (saut
de la question 2.02), surtout en milieu rural où l'accès à pied y est probable. Le surpeuplement suit la
définition ANSD (« plus de trois personnes par pièce ») opérationnalisée comme
au moins quatre personnes par pièce (ratio $\geq 4$), et non un ratio
strictement supérieur à 3 qui compterait à tort 3,5~pers./pièce comme
surpeuplé~; bien que l'EHCVM ne distingue pas les « pièces pour dormir » de la
définition-type (nombre total de pièces occupées), cette lecture reproduit
exactement les taux ANSD (13,3~/~12,7~/~11,4). Le NEET des 15-17 ans est,
lui, \emph{écarté} de l'agrégat. Sa valeur ANSD (85,7\,\%) se reproduit
exactement, mais elle correspond à la seule \emph{absence d'emploi}
(variable d'activité des 7 derniers jours différente de « occupé »), sans la
condition de non-scolarisation pourtant inscrite dans sa définition (« ne
poursuivant pas d'études ni de formation »)~: elle compte donc comme privés
les élèves — la quasi-totalité des 15-17 ans n'étant pas en emploi —, ce qui
la rend inutilisable comme indicateur de privation éducative. Un NEET
correctement défini (sans emploi \emph{et} non scolarisé) vaut 33,7\,\%, mais
n'est pas publié par l'ANSD~; conformément au principe de ne retenir que les
indicateurs validables, la dimension Éducation des 15-17 ans repose sur le seul
illettrisme (notre NEET correct, 33,7\,\%, reste rapporté à titre descriptif).
La séparation parentale (« enfant ne vivant pas avec ses deux parents
biologiques »), troisième indicateur de la dimension Protection, reproduit
étroitement les taux publiés par l'ANSD~: 36,5~/~43,0~/~52,8 contre
36,5~/~41,8~/~52,8, les groupes 0-4 et 15-17 coïncidant exactement.

Au-delà de la comparaison indicateur par indicateur, le tableau~\ref{tab:validation_globale}
confronte l'incidence N-MODA globale (nationale) obtenue dans ce mémoire à
celle publiée par l'ANSD.

\begin{table}[H]
  \centering
  \caption{Comparaison de l'incidence N-MODA globale : ce mémoire vs ANSD (EHCVM 2018, enfants 0-17 ans)}
  \label{tab:validation_globale}
  \zebre
  \begin{tabular}{|l|c|c|}
    \hline
\entete Statistique & Ce mémoire & ANSD (MODA 2024) \\
    \hline
    Incidence $H$ (\%)         & 56,4  & 50,7  \\ \hline
    Intensité $A$ (\%)         & 69,5  & 72,2  \\ \hline
    Indice ajusté $M_0 = H \times A$ & 0,392 & 0,366 \\
    \hline
  \end{tabular}

  {\footnotesize \textit{Source :} Calcul de l'auteur, EHCVM~; ANSD/UNICEF (MODA 2024) pour la colonne de comparaison.}
\end{table}

L'écart d'incidence (56,4\,\% contre 50,7\,\%) s'explique par deux facteurs.
D'une part, la nutrition est opérationnalisée de façon plus étroite que la
définition officielle~: elle ne repose ici que sur la sécurité alimentaire
(module FIES), sans le volet diversité des repas retenu par l'ANSD. D'autre
part, la dimension Santé est ici opérationnalisée de façon plus large que
dans l'agrégat de l'ANSD~: nous retenons l'union de ses deux indicateurs
officiels (combustible solide OU absence d'accès à pied à une structure de
santé), là où l'ANSD écarte le combustible de son agrégat multidimensionnel
pour manque de variance (92\,\% de privation, sans pouvoir discriminant), si
bien que sa dimension Santé se réduit à l'accès aux soins. Conserver le
combustible dans l'union élève mécaniquement la prévalence de la dimension et,
partant, l'incidence globale. Le solde de l'écart tient à des différences de
traitement statistique de l'échantillon par rapport au calcul officiel.
L'intensité, elle, reste proche de la valeur officielle.
